% !PRINTED 487
\weberpage{487}
%<<CONT
from this
\[
f(\omega)^8f(2\omega)^8+f_1(\omega)^8f_2(2\omega)^8=8+\frac{[f(\omega)^8+f_1(\omega)^8]^2}{f_2(\omega)^4},
\]
which by \S\,34 can easily be brought into the form:
\[
[f(\omega)^4f(2\omega)^4+f_1(\omega)^4f_2(2\omega)^4]^2=16+f_2(\omega)^{12}+\frac{64}{f_2(\omega)^{12}},
\]
and from this the root can be extracted:
\weq{5}{f(\omega)^4f(2\omega)^4+f_1(\omega)^4f_2(2\omega)^4=f_2(\omega)^6+\dfrac{8}{f_2(\omega)^6}.}
Now let $\omega$ satisfy the quadratic equation:
\weq{6}{2\,\omega^2+2\,r\,\omega+n=0}
with odd $n$, hence:
\weq{7}{\omega=\dfrac{-r+\sqrt{-m}}{2},\quad m=2\,n-r^2.}
If then
\weq{8}{\sqrt2\,x=f(\sqrt{-m})^2,\quad \text{or}\ =f_1(\sqrt{-m})^2}
% [sic] German source omits $\sqrt2\,x$ before the second equality
is put, according as $r$ is odd or even, then there results
\weq{9}{f_2(\omega)^2=e^{-\frac{r\pi i}{12}}\dfrac{\sqrt2}{x}.}
Furthermore, from
\[
2\,\omega=-\frac{n}{\omega}-2\,r
\]
there follows, by the formulae of \S\,34:
\[
f(2\omega)=f\!\left(\frac{\omega}{n}\right)e^{\frac{r\pi i}{12}}
\]
\weq{10}{\begin{aligned}
f_1(2\omega)&=f_2\!\left(\dfrac{\omega}{n}\right)e^{\frac{r\pi i}{12}} \\
f_2(2\omega)&=f_1\!\left(\dfrac{\omega}{n}\right)e^{-\frac{r\pi i}{6}},
\end{aligned}}
so that equation (5) gives:
\weq{11}{f(\omega)^4f\!\left(\dfrac{\omega}{n}\right)^4+(-1)^r f_1(\omega)^4f_1\!\left(\dfrac{\omega}{n}\right)^4=e^{-\frac{r\pi i}{12}}\sqrt8\left(x^3+\dfrac{e^{-\frac{r\pi i}{2}}}{x^3}\right).}
